\documentclass[11pt,letterpaper]{article}

%%% IDIOMA
\usepackage[english,spanish]{babel}
\addto\shorthandsspanish{\spanishdeactivate{~<>}}

%%% CODIFICACIÓN Y FUENTES
\usepackage[utf8]{inputenc}
\usepackage[T1]{fontenc}
\usepackage{lmodern}
\usepackage{microtype}

%%% PAQUETES
\usepackage{graphicx}
\usepackage{amsmath}
\usepackage{amssymb}
\usepackage{booktabs}
\usepackage{array}
\usepackage{longtable}
\usepackage{multirow}
\usepackage{siunitx}
\usepackage{xcolor}
\usepackage{tcolorbox}
\usepackage{fancyhdr}
\usepackage[hidelinks]{hyperref}
\usepackage[margin=2.5cm]{geometry}
\usepackage{titlesec}
\usepackage{enumitem}
\usepackage{listings}
\usepackage{float}

%%% COLORES INSTITUCIONALES USM
\definecolor{usmblue}{RGB}{0,45,98}
\definecolor{usmgray}{RGB}{88,89,91}
\definecolor{usmlight}{RGB}{234,240,247}

%%% CONFIGURACIÓN DE LISTINGS PARA SQL
\lstset{
    language=SQL,
    basicstyle=\ttfamily\small,
    keywordstyle=\color{usmblue}\bfseries,
    stringstyle=\color{red!70!black},
    commentstyle=\color{usmgray}\itshape,
    backgroundcolor=\color{usmlight},
    frame=single,
    rulecolor=\color{usmblue},
    breaklines=true,
    keepspaces=true,
    showstringspaces=false,
    belowskip=1em,
    aboveskip=1em,
    inputencoding=utf8,
    extendedchars=true,
    literate={á}{{\'a}}1 {é}{{\'e}}1 {í}{{\'i}}1 {ó}{{\'o}}1 {ú}{{\'u}}1
             {Á}{{\'A}}1 {É}{{\'E}}1 {Í}{{\'I}}1 {Ó}{{\'O}}1 {Ú}{{\'U}}1
             {ñ}{{\~n}}1 {Ñ}{{\~N}}1 {ü}{{\"u}}1 {Ü}{{\"U}}1
}

%%% siunitx: formato chileno (punto de miles, coma decimal)
\sisetup{output-decimal-marker={,},group-separator={.},group-minimum-digits=4}
\newcommand{\clp}[1]{\$\num{#1}}

%%% TÍTULOS
\titleformat{\section}{\color{usmblue}\Large\bfseries}{\thesection.}{0.5em}{}
\titleformat{\subsection}{\color{usmblue}\large\bfseries}{\thesubsection.}{0.5em}{}
\titleformat{\subsubsection}{\color{usmgray}\normalsize\bfseries}{\thesubsubsection.}{0.5em}{}

%%% ENCABEZADO / PIE
\pagestyle{fancy}
\fancyhf{}
\renewcommand{\headrulewidth}{0.8pt}
\renewcommand{\footrulewidth}{0.4pt}
\fancyhead[L]{\footnotesize\textcolor{usmblue}{\textbf{Universidad Técnica Federico Santa María}}\\ \footnotesize\textcolor{usmgray}{ICN-292 --- Sistemas de Información para la Gestión}}
\fancyhead[R]{\footnotesize\textcolor{usmblue}{\textbf{Departamento de Industrias}}\\ \footnotesize\textcolor{usmgray}{Segundo Semestre 2026}}
\fancyfoot[C]{\footnotesize\textcolor{usmgray}{\thepage}}

%%% CAJAS
\newtcolorbox{cajaformula}{colback=usmlight,colframe=usmblue,boxrule=0.6pt,arc=2pt,left=6pt,right=6pt,top=4pt,bottom=4pt}
\newtcolorbox{cajanota}{colback=white,colframe=usmgray,boxrule=0.6pt,arc=2pt,left=6pt,right=6pt,top=4pt,bottom=4pt}

\begin{document}

%%--------------------------------------------------
%% PORTADA
%%--------------------------------------------------
\begin{titlepage}
\centering
\vspace*{1cm}
\IfFileExists{logo-usm.png}{\includegraphics[width=4.5cm]{logo-usm.png}\\[0.6cm]}{\vspace{0.6cm}}
{\color{usmblue}\rule{\textwidth}{2pt}}\\[0.4cm]
{\LARGE\bfseries\color{usmblue} Universidad Técnica Federico Santa María}\\[0.2cm]
{\large\color{usmgray} Departamento de Industrias}\\[0.2cm]
{\large\color{usmgray} ICN-292 --- Sistemas de Información para la Gestión}\\[0.3cm]
{\color{usmblue}\rule{\textwidth}{2pt}}\\[2.5cm]

{\Huge\bfseries Laboratorio 5: SQL y Bases de Datos\\[0.3cm] Creación de la Base y Consultas}\\[0.6cm]
{\Large\itshape Andes Bike}\\[0.4cm]
{\large Creación del modelo relacional, carga oficial y diez consultas de validación}\\[2.5cm]

\begin{tabular}{rl}
\textbf{Autor:} & Matías Gómez \\[0.3em]
\textbf{Rol USM:} & 202360624-3 \\[0.3em]
\textbf{Paralelo:} & 100 \\[0.3em]
\textbf{Profesores:} & José Miguel González Paul \\[0.3em]
 & José Luis Sáez Tamayo \\[0.3em]
 & Patricio Cifuentes \\[0.3em]
\textbf{Repositorio GitHub:} & \url{https://github.com/Maty0401/ICN292-Lab5-Gomez-Matias} \\[0.3em]
\textbf{Fecha:} & 11 de noviembre de 2026 \\
\end{tabular}

\vfill
{\color{usmblue}\rule{\textwidth}{1pt}}
\end{titlepage}

\newpage
%%--------------------------------------------------
\section*{Consultas SQL}
\addcontentsline{toc}{section}{Consultas SQL}
%%--------------------------------------------------

A continuación se presenta, para cada una de las diez consultas solicitadas, el código SQL tal como fue ejecutado y la captura del resultado devuelto por MySQL Workbench sobre la carga oficial de la base \texttt{andesbike}.


\subsection*{Consulta 1}
Liste todos los vendedores con la cantidad de ventas que registraron, incluidos los que no registraron ninguna. Ordene de mayor a menor cantidad y, a igualdad, por nombre. Los que no tienen ventas deben mostrar cero, no quedar en blanco.

\begin{lstlisting}
SELECT  v.nombre                AS vendedor,
        COUNT(ve.id_venta)      AS cantidad_ventas
FROM       vendedores v
LEFT JOIN  ventas     ve  ON ve.id_vendedor = v.id_vendedor
GROUP BY   v.id_vendedor, v.nombre
ORDER BY   cantidad_ventas DESC, v.nombre ASC;
\end{lstlisting}
\begin{figure}[H]
    \centering
    \includegraphics[width=0.9\textwidth]{Captura de pantalla 2026-11-11 143407.png}
    \caption{Resultado de la Consulta 1 en MySQL Workbench.}
\end{figure}

\subsection*{Consulta 2}
Para cada tienda muestre su nombre, cuántas ventas se registraron y cuántas unidades se vendieron, incluidas las tiendas que no tienen ningún vendedor. Ordene por unidades de mayor a menor.

\begin{lstlisting}
SELECT  t.nombre                          AS tienda,
        COUNT(ve.id_venta)                AS ventas_registradas,
        COALESCE(SUM(ve.unidades), 0)     AS unidades_vendidas
FROM       tiendas    t
LEFT JOIN  vendedores v   ON v.id_tienda    = t.id_tienda
LEFT JOIN  ventas     ve  ON ve.id_vendedor = v.id_vendedor
GROUP BY   t.id_tienda, t.nombre
ORDER BY   unidades_vendidas DESC;
\end{lstlisting}
\begin{figure}[H]
    \centering
    \includegraphics[width=0.9\textwidth]{Captura de pantalla 2026-11-11 143431.png}
    \caption{Resultado de la Consulta 2 en MySQL Workbench.}
\end{figure}

\subsection*{Consulta 3}
Liste las tiendas que no tienen ningún vendedor asignado.

\begin{lstlisting}
SELECT  t.id_tienda,
        t.nombre   AS tienda,
        t.comuna
FROM       tiendas    t
LEFT JOIN  vendedores v  ON v.id_tienda = t.id_tienda
WHERE      v.id_vendedor IS NULL
ORDER BY   t.nombre;
\end{lstlisting}
\begin{figure}[H]
    \centering
    \includegraphics[width=0.9\textwidth]{Captura de pantalla 2026-11-11 143443.png}
    \caption{Resultado de la Consulta 3 en MySQL Workbench.}
\end{figure}

\subsection*{Consulta 4}
Liste los modelos que no se vendieron ni una sola vez en junio de 2025.

\begin{lstlisting}
SELECT  b.modelo,
        b.categoria,
        b.precio
FROM       bicicletas b
LEFT JOIN  ventas     ve
       ON  ve.id_bicicleta = b.id_bicicleta
      AND  ve.fecha >= '2025-06-01'
      AND  ve.fecha  < '2025-07-01'
WHERE      ve.id_venta IS NULL
ORDER BY   b.modelo;
\end{lstlisting}
\begin{figure}[H]
    \centering
    \includegraphics[width=0.9\textwidth]{Captura de pantalla 2026-11-11 143516.png}
    \caption{Resultado de la Consulta 4 en MySQL Workbench.}
\end{figure}

\begin{cajanota}
\textbf{Ubicación de la condición del mes:} la condición de junio de 2025 está puesta en el \texttt{ON} del \texttt{LEFT JOIN}, de modo que el join trae únicamente las ventas de ese mes y el \texttt{WHERE ve.id\_venta IS NULL} deja los modelos que no tuvieron ninguna; si esa misma condición se hubiera puesto en el \texttt{WHERE}, se evaluaría después del join sobre filas ya rellenadas con \texttt{NULL}, chocaría con la condición \texttt{IS NULL} y la consulta habría devuelto cero filas.
\end{cajanota}

\subsection*{Consulta 5}
Liste los vendedores que no han registrado ninguna venta, mostrando también el nombre de su tienda. Ninguno de ellos puede quedar fuera del resultado.

\begin{lstlisting}
SELECT  v.nombre                                      AS vendedor,
        COALESCE(t.nombre, '(sin tienda asignada)')   AS tienda,
        v.fecha_ingreso
FROM       vendedores v
LEFT JOIN  tiendas    t   ON t.id_tienda    = v.id_tienda
LEFT JOIN  ventas     ve  ON ve.id_vendedor = v.id_vendedor
WHERE      ve.id_venta IS NULL
ORDER BY   v.nombre;
\end{lstlisting}
\begin{figure}[H]
    \centering
    \includegraphics[width=0.9\textwidth]{Captura de pantalla 2026-11-11 143529.png}
    \caption{Resultado de la Consulta 5 en MySQL Workbench.}
\end{figure}

\subsection*{Consulta 6}
Liste los pares de vendedores que trabajan en la misma tienda, poniendo primero al que ingresó antes. Muestre los dos nombres, las dos fechas de ingreso y la tienda. Cada par debe aparecer una sola vez.

\begin{lstlisting}
SELECT  t.nombre          AS tienda,
        a.nombre          AS vendedor_mas_antiguo,
        a.fecha_ingreso   AS ingreso_mas_antiguo,
        b.nombre          AS vendedor_mas_nuevo,
        b.fecha_ingreso   AS ingreso_mas_nuevo
FROM  vendedores a
JOIN  vendedores b
   ON  a.id_tienda = b.id_tienda
  AND (a.fecha_ingreso <  b.fecha_ingreso
   OR (a.fecha_ingreso =  b.fecha_ingreso
       AND a.id_vendedor < b.id_vendedor))
JOIN  tiendas t ON t.id_tienda = a.id_tienda
ORDER BY t.nombre, a.fecha_ingreso, b.fecha_ingreso;
\end{lstlisting}
\begin{figure}[H]
    \centering
    \includegraphics[width=0.9\textwidth]{Captura de pantalla 2026-11-11 143621.png}
    \caption{Resultado de la Consulta 6 en MySQL Workbench.}
\end{figure}

\subsection*{Consulta 7}
Para cada tienda, calcule el monto total vendido en bicicletas de las categorías \texttt{'Montaña'} y \texttt{'Ruta'}, dejando solo las tiendas que superan los \clp{2000000}. Ordene de mayor a menor.

\begin{lstlisting}
SELECT  t.nombre                          AS tienda,
        SUM(ve.unidades * b.precio)       AS monto_total
FROM  tiendas    t
JOIN  vendedores v   ON v.id_tienda     = t.id_tienda
JOIN  ventas     ve  ON ve.id_vendedor  = v.id_vendedor
JOIN  bicicletas b   ON b.id_bicicleta  = ve.id_bicicleta
WHERE     b.categoria IN ('Montaña', 'Ruta')
GROUP BY  t.id_tienda, t.nombre
HAVING    monto_total > 2000000
ORDER BY  monto_total DESC;
\end{lstlisting}
\begin{figure}[H]
    \centering
    \includegraphics[width=0.9\textwidth]{Captura de pantalla 2026-11-11 143648.png}
    \caption{Resultado de la Consulta 7 en MySQL Workbench.}
\end{figure}

\subsection*{Consulta 8}
Liste el modelo y el precio de las bicicletas cuyo precio supera el promedio del catálogo, de mayor a menor.

\begin{lstlisting}
SELECT  b.modelo,
        b.precio
FROM    bicicletas b
WHERE   b.precio > (SELECT AVG(precio) FROM bicicletas)
ORDER BY b.precio DESC;
\end{lstlisting}
\begin{figure}[H]
    \centering
    \includegraphics[width=0.9\textwidth]{Captura de pantalla 2026-11-11 143700.png}
    \caption{Resultado de la Consulta 8 en MySQL Workbench.}
\end{figure}

\begin{cajanota}
\textbf{Valor del promedio y su origen:} el promedio del catálogo es \clp{657490.00} y se obtuvo con la subconsulta escalar \texttt{SELECT AVG(precio) FROM bicicletas}, es decir, promediando los ocho modelos del catálogo (\num{5259920} / 8 = \num{657490}) y no las unidades efectivamente vendidas.
\end{cajanota}

\subsection*{Consulta 9}
Liste los vendedores cuyo promedio de unidades por venta supera el promedio de unidades de todas las ventas de la cadena. Muestre el nombre y su promedio, de mayor a menor y, a igualdad, por nombre.

\begin{lstlisting}
SELECT  v.nombre            AS vendedor,
        AVG(ve.unidades)    AS promedio_unidades
FROM  vendedores v
JOIN  ventas     ve  ON ve.id_vendedor = v.id_vendedor
GROUP BY  v.id_vendedor, v.nombre
HAVING    AVG(ve.unidades) > (SELECT AVG(unidades) FROM ventas)
ORDER BY  promedio_unidades DESC, v.nombre ASC;
\end{lstlisting}
\begin{figure}[H]
    \centering
    \includegraphics[width=0.9\textwidth]{Captura de pantalla 2026-11-11 143712.png}
    \caption{Resultado de la Consulta 9 en MySQL Workbench.}
\end{figure}

\subsection*{Consulta 10}
Liste todos los vendedores con el monto total que vendieron, incluidos los que no vendieron nada, de mayor a menor monto.

\begin{lstlisting}
SELECT  v.nombre                       AS vendedor,
        SUM(ve.unidades * b.precio)    AS monto_total
FROM       vendedores v
LEFT JOIN  ventas     ve  ON ve.id_vendedor  = v.id_vendedor
LEFT JOIN  bicicletas b   ON b.id_bicicleta  = ve.id_bicicleta
GROUP BY   v.id_vendedor, v.nombre
ORDER BY   monto_total DESC;
\end{lstlisting}
\begin{figure}[H]
    \centering
    \includegraphics[width=0.9\textwidth]{Captura de pantalla 2026-11-11 143730.png}
    \caption{Resultado de la Consulta 10 en MySQL Workbench.}
\end{figure}
\begin{figure}[H]
    \centering
    \includegraphics[width=0.9\textwidth]{Captura de pantalla 2026-11-11 143738.png}
    \caption{Resultado de la Consulta 10 en MySQL Workbench (continuación).}
\end{figure}

\begin{cajanota}
\textbf{Por qué aparecen \texttt{NULL} y no cero:} Fernanda Lillo y Rubén Castro no muestran cero sino \texttt{NULL} porque el \texttt{LEFT JOIN} sí los conserva, pero no tienen ninguna fila de venta que sumar y \texttt{SUM()} sobre un conjunto vacío devuelve \texttt{NULL}; para que mostraran cero habría que agregar \texttt{COALESCE} a la suma, escribiendo \texttt{COALESCE(SUM(ve.unidades * b.precio), 0) AS monto\_total}.
\end{cajanota}

\end{document}
